\chapter{Training Configuration}
\label{app:training_config}

\section{Optimiser and Learning Rate Schedule}
\label{app:sec:optim}

\textbf{AdamW} \quad At each step, the bias-corrected first and second moment
estimates $\hat{m}_t$ and $\hat{v}_t$ are computed with decay coefficients
$\beta_1$ and $\beta_2$, and the parameter update is:
\[
    \theta_t = \theta_{t-1}
        - \eta\,\frac{\hat{m}_t}{\sqrt{\hat{v}_t} + \varepsilon}
        - \eta\lambda\,\theta_{t-1},
\]
where $\eta$ is the learning rate, $\lambda$ the weight-decay coefficient, and
$\varepsilon$ a small constant for numerical stability.

\textbf{Cosine annealing} \quad The learning rate at step $t$ of a schedule of
total length $T$ is:
\[
    \eta_t = \eta_{\min}
    + \tfrac{1}{2}(\eta_{\max} - \eta_{\min})
    \!\left(1 + \cos\!\left(\frac{t}{T}\pi\right)\right).
\]
A known risk of this schedule is entrapment in a local minimum if the learning
rate decays below the threshold needed to escape it before convergence is
reached.

\section{Model Dimension and Parameter Count}
\label{app:model-parameter-count}

This section derives the trainable parameter count of the CSDI-based diffusion
backbone for the two experimental configurations, which differ only in the
number of target features~$K$.

\subsection{Notation and Hyperparameters}
Here is some notation: \(K\) is \texttt{target\_dim}, 
\(E_f\) is \texttt{featureemb},
\(E_t\) is \texttt{timeemb}, \(C\) is \texttt{channels}, \(D\) is \texttt{diffusion\_embedding\_dim}, \(R\) is \texttt{num\_layers}.

The side-information width per (feature, time) location is
$S = E_t + E_f + 1$,
where the $+1$ is the conditioning mask channel.
This is independent of $K$ because $E_f$ is the output dimension of the
feature embedding lookup, not the vocabulary size.
The Transformer feed-forward width is $d_{\mathrm{ff}} = 4C$.
Fixed values:
\[
E_f = 64,\quad
E_t = 128,\quad
C = 128,\quad
D = 256,\quad
S = 193,\quad
I = 2,\quad
R = 6.
\]

\subsection{Shared Parameters}

\textbf{Non-backbone}
\begin{itemize}
    \item \textbf{Time-embedding projection}: $\mathrm{Linear}(E_t, E_t)$: \(E_t^2 + E_t = 16{,}512.\)
    \item \textbf{Outside residual blocks (\texttt{diff\_CSDI})}: input projection $\mathrm{Conv1d}(I,C,1)$, output projection $\mathrm{Conv1d}(C,1,1)$, and the two-layer diffusion embedding $2\,\mathrm{Linear}(D,D)$
\end{itemize}
The total count of Non-backbone element is:
\[
N_{\mathrm{outside}}
= (IC+C) + (C+1) + 2(D^2+D)
= 384 + 129 + 131{,}584
= 132{,}097.
\]

\textbf{Residual Block (\texttt{ResidualBlock})}
\begin{itemize}
    \item \textbf{Four projection layers}: diffusion $\mathrm{Linear}(D,C)$, conditioning $\mathrm{Conv1d}(S,2C,1)$, mid and output $\mathrm{Conv1d}(C,2C,1)$ for each residual block:
\[
N_{\mathrm{proj}}
= DC + 2CS + 4C^2 + 7C
= 148{,}608.
\]
    \item \textbf{one Transformer encoder layer}: $d_{\mathrm{model}}=C$, $d_{\mathrm{ff}}=4C$, contributions from multi-head attention $4C^2+4C$, feed-forward network $8C^2+5C$, and two layer-norm layers $4C$ for each residual block:
\[
N_{\mathrm{trans}}
= 12C^2 + 13C
= 198{,}272.
\]
\end{itemize}

\subsection{K-dependent Parameters}

Two components depend on $K$.
\begin{itemize}
    \item \textbf{Feature embedding}: $\mathrm{Embedding}(K,E_f)$: contributes $K\cdot E_f$ parameters.
    \item \textbf{Feature Transformer}: when $K>1$, each residual block contains a second Transformer encoder layer of identical architecture to the time Transformer, operating across the feature axis ($N_{\mathrm{trans}}$ additional parameters per block). For $K=1$ this layer is instantiated, but never updated. The reason is just code compatability across phases. 
\end{itemize}

\begin{table}[h]
\centering
\begin{tabular}{lcc}
\toprule
Component & $K=1$ & $K=4$ \\
\midrule
Feature embedding $K E_f$ & $64$ & $256$ \\
Feature Transformer per block & $0$ & $198{,}272$ \\
\midrule
Parameters per residual block $N_{\mathrm{res}}$ & $346{,}880$ & $545{,}152$ \\
\bottomrule
\end{tabular}
\end{table}

\subsection{Total Parameter Count}

\[
N_{\mathrm{total}}
= \underbrace{\bigl(K E_f + E_t^2 + E_t\bigr)}_{\text{non-backbone}}
+ \underbrace{N_{\mathrm{outside}}}_{\text{outside blocks}}
+ R\cdot N_{\mathrm{res}}.
\]

\begin{table}[h]
\centering
\begin{tabular}{lcc}
\toprule
 & $K=1$ & $K=4$ \\
\midrule
Non-backbone & $16{,}576$ & $16{,}768$ \\
Outside residual blocks & \multicolumn{2}{c}{$132{,}097$} \\
$R \cdot N_{\mathrm{res}}$ & $6\times 346{,}880 = 2{,}081{,}280$ & $6\times 545{,}152 = 3{,}270{,}912$ \\
\midrule
\textbf{Total} & $\mathbf{2{,}229{,}953}$ & $\mathbf{3{,}419{,}777}$ \\
\bottomrule
\end{tabular}
\end{table}

The two configurations differ by $1{,}189{,}824$ parameters
($6 \times N_{\mathrm{trans}} + (K_4 - K_1)\cdot E_f
= 1{,}189{,}632 + 192$),
entirely attributable to the six feature Transformer layers absent in the
univariate case and the larger feature embedding vocabulary.

\section{Attention Memory and Sequence Length Constraint}
\label{app:sec:attention_memory}

The CSDI self-attention computes a score matrix of shape \((B \times K,\, H,\, L,\, L)\), and
PyTorch selects among three SDPA kernels at runtime. Flash SDPA avoids materialising
the full matrix but requires more CUDA cores than the available MIG slices provide.
Memory-Efficient SDPA encounters an internal limit of 65,535 on the
\(\text{batch} \times \text{heads}\) dimension: with \(K = 4\), \(B = 64\), and
\(L = 2048\), this product reaches 524,288. The only remaining fallback, Math SDPA,
materialises the full \(L \times L\) attention matrix, costing approximately 17 GB for
attention alone at \(L = 2048\), which exceeds the 40 GB budget once model parameters
and activations are included. Reducing to \(L = 512\) brings the attention cost within
budget while still covering approximately two years of daily returns; the stride of 100
maintains the same ratio of overlap as \(2048/400\).